Data management is essential for the continuity of research, outreach, and administrative activities in universities. However, datasets produced in these institutions often remain scattered across local files and unintegrated systems, lacking documentation, access control, and history preservation, which hinders their discovery, interpretation, and reuse. Addressing this problem, this work describes the development of DataRural, an institutional web platform created to centralize and organize the publication and sharing of datasets at the Universidade Federal Rural do Rio de Janeiro (UFRRJ). The solution enables uploading CSV datasets accompanied by descriptive metadata, versioning with preservation of previous states, visibility control between public and private, and collaboration through groups with distinct roles. The platform also offers tabular previews of files and individual or bundled downloads. This work demonstrates the feasibility of combining publication, traceability, and controlled access to academic data within an institutional environment, and suggests avenues for future enhancements, such as support for new formats, the provision of an API for programmatic consumption, and usability evaluation with the university community.

\vspace{0.5cm}
\noindent\textbf{Keywords:} data management, institutional repository, datasets, web platform, versioning, UFRRJ.
